%--- iliad ---
% title: Reward Learning Theory
% summary: >-
%   RLHF grades a policy by a proxy for the true reward: what the evaluator can
%   see. We derive by hand, in a toy example with partial observability, how a
%   policy can hide its failures from the evaluator and inflate the evaluator's
%   perception of its return.
% slides:
%   url: https://drive.google.com/drive/folders/18NswW6zYnDp9KLKjG_1BwctSPITOjZ-A
%   title: Intro to Reward Learning Theory
%--- end ---
\documentclass[11pt]{article}

\usepackage{mathtools}
\usepackage{float}
\usepackage[numbers]{natbib}

% iliad.sty: local copy first (Overleaf), else the shared ../iliad.sty. It
% supplies the definition/theorem/fact/remark family, exercise + solution,
% callout, hint, hyperref + cleveref — everything the upstream preamble built
% by hand (aliascnt, tcolorbox, commenting.sty).
\IfFileExists{iliad.sty}{\usepackage[boxes]{iliad}}{\usepackage[boxes]{../iliad}}

\hypersetup{colorlinks=true, linkcolor=blue, urlcolor=blue, citecolor=blue}

% Math macros — as upstream.
\newcommand{\cS}{\mathcal{S}}
\newcommand{\cA}{\mathcal{A}}
\newcommand{\cO}{\mathcal{O}}
\newcommand{\cB}{\mathcal{B}}
\newcommand{\cT}{\mathcal{T}}
\newcommand{\E}{\mathbb{E}}
\newcommand{\R}{\mathbb{R}}
\newcommand{\Gobs}{G_{\mathrm{obs}}}
\newcommand{\Jobs}{J_{\mathrm{obs}}}
\newcommand{\Erro}{E^{+}}
\newcommand{\Erru}{E^{-}}
\DeclareMathOperator*{\argmax}{arg\,max}

\title{Reward Learning Theory}
\author{
  \authorname{Leon Lang}\\
  \affiliation{Iliad}
  \and
  \authorname{Joar Skalse}\\
  \affiliation{Deducto Limited, King’s College London}
}
\date{\today}

\begin{document}
\maketitle

%===============================================================================
% Front matter from the A.4 tab of the Iliad Mega Doc, transcribed verbatim.
\section*{Prerequisites}

It is useful to know the basics of reinforcement learning, reinforcement learning from human feedback (RLHF), and AI alignment, as taught in other modules of this course.

\begin{learningoutcomes}
  \begin{itemize}
    \item learn the place of reward learning in the AI alignment landscape;
    \item can reason about the (speculative) assumptions that motivate this research area as a whole;
    \item understand the strong conditions under which the simplest reward learning technique, RLHF via trajectory comparisons, would lead to an aligned objective;
    \item can reason about situations where such conditions do not hold, like underspecification and misspecification, and efforts to account for, or correct, such issues;
    \item and understand how reward learning can be embedded into the framework of assistance games, which can be regarded as one conceptualization of the entire alignment problem.
  \end{itemize}
\end{learningoutcomes}

\section*{Roadmap for the day (as taught in the April 2026 Iliad Intensive)}

\begin{itemize}
  \item 10:00 – 11:30 Lecture and questions
  \item 11:30 – 12:30: Exercise block 1
  \item 12:30 – 13:30: Lunch break
  \item 13:30 – 14:30: Exercise block 2
  \item 14:30 – 15:30: Guest lecture by Richard Ngo
  \item 15:30 – 16:00: Break
  \item 16:00 – 17:30: Guest lecture by Joar Skalse
  \item 17:30 – 18:30: Discussions, feedback, and wrap-up
\end{itemize}

\section*{Content}

\subsection*{Fast-track}

Simply read the lecture slides of the first lecture. If you have more time, do the first exercise on partial observability as one instantiation of misspecification.

\subsection*{Main content}

\begin{itemize}
  \item \href{https://docs.google.com/presentation/d/1cG2M68gj8osmrse97bza9cgCqqwoFXBvDIGpkv-HVgA/edit?usp=drive_link}{Lecture slides: Intro to Reward Learning Theory by Leon Lang}
  \item Reward Learning Theory: Exercises (also available without solutions). The exercise sheet is based on the following two papers:
  \begin{itemize}
    \item \href{https://arxiv.org/abs/2402.17747}{When Your AIs Deceive You: Challenges of Partial Observability in Reinforcement Learning from Human Feedback}
    \item \href{https://people.eecs.berkeley.edu/~russell/papers/neurips20ws-assistance}{Benefits of Assistance over Reward Learning}
  \end{itemize}
  \item \href{https://drive.google.com/file/d/1WLC-HILvYGmXo9kqUGEvekK6scURXQBl/view?usp=sharing}{Lecture slides: Towards a Formal Theory of Reward Learning, With Application to Inverse Reinforcement Learning} by Joar Skalse. More context:
  \begin{itemize}
    \item \href{https://www.lesswrong.com/s/TEybbkyHpMEB2HTv3/p/pJ3mDD7LfEwp3s5vG}{The Theoretical Reward Learning Research Agenda: Introduction and Motivation} by Joar Skalse
  \end{itemize}
\end{itemize}

%===============================================================================
% The exercise sheet (Leon Lang, "Reward Learning Theory: Exercises"), as it
% compiles upstream. Its Setup stays unnumbered so "Hiding failures" is section
% 1 and the problems keep their upstream numbers, 1.1–1.10.




%===================================

\section*{Setup}

This sheet (\Cref{sec:hiding}) works through one
concrete example from~\citet{Lang2024PORLHF} to see how naive RLHF can
fail when the human evaluator only \emph{partially} observes the
environment they are providing feedback on. By the end of the sheet you
will have derived, by hand, the conditions under which an RLHF-optimal
policy hides information from the human in a way that systematically
inflates the human's perception of the policy's return --- the failure
mode the paper calls \textbf{deceptive inflation}. We close with an
informal discussion of the dual failure mode, \textbf{overjustification},
in which the agent pays real reward to make its (already-good) behavior
look as good as it is, and briefly discuss how such failure modes
motivate approaches like AI safety via debate.

% Part 2 intro; re-enable together with \include{sec_embedding} below.
% \textbf{Part~2 (\Cref{sec:embedding}: Reward learning embeds into
% assistance)} steps back from RLHF specifically and asks whether the
% reward-learning paradigm is itself the most general way to learn from
% human feedback. Following~\citet{Shah2020Assistance}, we introduce the
% \textbf{assistance} framework, in which the human is part of the
% environment and the true reward is a hidden parameter the agent must
% infer. The big exercise of this part proves that every active reward
% learning problem reduces to a particular --- restricted --- kind of
% assistance problem, making precise in what sense reward learning is a
% special case of assistance.

We assume familiarity with finite-horizon Markov decision processes;
the additional structure needed --- observation kernels, human beliefs,
the observation return $\Gobs$ and observation value $\Jobs$ --- is
introduced in highlighted boxes as we go.

% ================================================

\begin{definition}[MDP and trajectories]\label{def:mdp}
A finite-horizon \textbf{Markov decision process} is a tuple
$(\cS, \cA, \cT, R, P_0, T, \gamma)$ with finite state space $\cS$, finite
action space $\cA$, transition kernel
$\cT : \cS \times \cA \to \Delta(\cS)$,
reward function $R : \cS \to \R$, initial-state
distribution $P_0 \in \Delta(\cS)$, horizon $T \in \mathbb{N}$, and discount
$\gamma \in (0,1]$. A \textbf{state trajectory} is a sequence
$\vec s = s_0 s_1 \cdots s_T$, and its \textbf{return} is
$$
G(\vec s) \;=\; \sum_{t=0}^{T} \gamma^{t}\, R(s_t).
$$
A \textbf{policy} $\pi : \cS \to \Delta(\cA)$ induces a distribution $P^\pi$
over state trajectories. The \textbf{policy evaluation function} $J$
assigns to each policy its expected return,
$$
J(\pi) \;=\; \E_{\vec s \sim P^\pi}[G(\vec s)],
$$
which we will also call the policy's \textbf{true value} (to contrast
later with the value the human \emph{thinks} it is evaluating).
\end{definition}

% ================================================

\section{Hiding failures}\label{sec:hiding}

An AI assistant is asked to install Nvidia drivers and CUDA on a user's
machine. It can attempt the CUDA install with default logging
(action~$a_C$, ``C'' for CUDA), or it can append \texttt{/dev/null} to the command
(action~$a_H$, ``H'' for hide), which suppresses any error message if the
installation fails. The user wants CUDA installed but \emph{also} dislikes
having errors silently hidden.

\paragraph{The MDP.} The full MDP is depicted in
\Cref{fig:cuda_mdp}. We use a finite horizon $T = 3$ and $\gamma = 1$:

\begin{figure}[H]
    \centering
    \includegraphics[width=\linewidth]{fig/expanded_example_figure1.pdf}
    \caption{The CUDA-installation MDP and its observation kernel
    \citep[Figure~6A]{Lang2024PORLHF}. Edges are labeled
    by the action triggering the transition; the small number in the top
    right of each state box is its reward; the small symbol in the
    bottom-right of each state box is its observation under $O$.}
    \label{fig:cuda_mdp}
\end{figure}

\begin{itemize}
    \item $\cS = \{S, I, W, W_H, L, L_H, T\}$. ($S$ = start; $I$ = drivers
        installed; $W$ = CUDA installed with default logging; $W_H$ = CUDA
        installed via the \texttt{/dev/null} trick; $L$ = CUDA install failed,
        error visible; $L_H$ = CUDA install failed, error hidden; $T$ =
        the absorbing, ``terminal'' MDP state in which the agent has stopped acting.)
    \item $\cA = \{a_I, a_C, a_H, a_T\}$. ($a_I$ installs drivers; $a_C$
        attempts CUDA install; $a_H$ attempts CUDA install with \texttt{/dev/null};
        $a_T$ stops acting, sending the MDP to $T$.)
    \item $P_0(S) = 1$, and from state $T$ every action goes to $T$.
    \item Transitions: $S \xrightarrow{a_I} I$. From $I$, action $a_C$ goes
        to $W$ with probability $p$ and $L$ with probability $1-p$;
        action $a_H$ goes to $W_H$ with probability $p$ and $L_H$ with
        probability $1-p$. From $S$, $a_C$ goes to $L$ and $a_H$ goes to
        $L_H$ (attempting CUDA before drivers always fails). Any other
        action transitions to $T$.
    \item Rewards: $R(S) = R(T) = 0$, $R(I) = 1$, $R(W) = R(W_H) = 10$,
        $R(L) = -5$, $R(L_H) = -5 - r$, where $r \ge 0$ is the user's
        penalty for hidden errors.
\end{itemize}

We restrict attention to deterministic policies. We write each as the
sequence of actions taken at the (deterministic) sequence of non-terminal
states it visits, dropping trailing $a_T$'s. Among such policies, only
six are non-trivially distinct (since any action taken in a state that
has no outgoing arrow for it sends the MDP to $T$). Of these six, the
four that we will analyze are
$$
[a_T], \quad [a_I a_T], \quad [a_I a_C a_T], \quad [a_I a_H a_T].
$$
The two omitted policies $[a_C a_T]$ and $[a_H a_T]$ attempt CUDA
\emph{before} installing drivers, which always fails; they are dominated
by $[a_T]$ in true value and add nothing of interest to the analysis
below.

% --- conceptual sanity check ---
\begin{exercise}
\label{prob:concept_setup}
(Conceptual.) In one sentence
each:
\begin{itemize}
    \item Why does the user reward $W$ and $W_H$ identically (both $+10$)?
    \item Why does the user reward $L_H$ strictly less than $L$ (i.e.\ why
        $r > 0$)?
\end{itemize}
\end{exercise}

\begin{solution}[prob:concept_setup]
$W$ and $W_H$ both correspond to a successful CUDA install. The user only
cares about the install succeeding, not about whether the agent
\emph{would have} hidden errors had it failed; on a successful run, no
error needed hiding. Thus the true reward is the same.

In contrast, $L_H$ is a failure where the agent has actively suppressed
the error message, depriving the user of information. The user prefers
to see the failure (state $L$) over having it hidden (state $L_H$);
the penalty $r > 0$ encodes this preference.
\end{solution}

% ----------------------------------------------------

\begin{exercise}
\label{prob:trueG}
For each of the eight state trajectories
listed below, compute the true return $G(\vec s)$:
$$
STTT,\; SL_HTT,\; SLTT,\; SITT,\; SIL_HT,\; SILT,\; SIWT,\; SIW_HT.
$$
\end{exercise}

\begin{solution}[prob:trueG]
By definition $G(\vec s) = \sum_{t=0}^{3} R(s_t)$. Plugging in the rewards:
\begin{center}
\begin{tabular}{c|cccccccc}
$\vec s$    & $STTT$ & $SL_HTT$ & $SLTT$ & $SITT$ & $SIL_HT$ & $SILT$ & $SIWT$ & $SIW_HT$ \\\hline
$G(\vec s)$ & $0$    & $-5-r$   & $-5$   & $1$    & $-4-r$   & $-4$   & $11$   & $11$
\end{tabular}
\end{center}
\end{solution}

% ----------------------------------------------------

\begin{definition}[Observation kernel]\label{def:obs}
An \textbf{observation kernel} on $\cS$ is a deterministic map
$O : \cS \to \Omega$ to a finite set $\Omega$ of observations. For a state
trajectory $\vec s = s_0 \cdots s_T$, we write
$\vec O(\vec s) = O(s_0) \cdots O(s_T)$ for the corresponding observation
trajectory. (More generally, one can take $O : \cS \to \Delta(\Omega)$ to
be stochastic, but this exercise sheet only needs the deterministic case.)
\end{definition}

For the CUDA example, the observation kernel is given by
\begin{center}
\begin{tabular}{c|ccccccc}
$s$ & $S$ & $I$ & $W$ & $W_H$ & $L$ & $L_H$ & $T$ \\\hline
$O(s)$ & $o_\emptyset$ & $o_I$ & $o_W$ & $o_W$ & $o_L$ & $o_\emptyset$ & $o_\emptyset$
\end{tabular}
\end{center}
The observations $o_I, o_W, o_L, o_\emptyset$ correspond, respectively, to a
log message confirming driver install, confirming CUDA install, reporting a
CUDA failure, and no log message at all.

% ----------------------------------------------------

\begin{exercise}
\label{prob:obsclasses}
Identify all pairs of trajectories from
\Cref{prob:trueG} that produce the \emph{same} observation trajectory under
$\vec O$. For each such pair, write down the shared observation trajectory.
\end{exercise}

\begin{solution}[prob:obsclasses]
Three pairs collide:
\begin{itemize}
    \item $STTT$ and $SL_HTT$ both produce $o_\emptyset o_\emptyset o_\emptyset o_\emptyset$ (empty log).
    \item $SITT$ and $SIL_HT$ both produce $o_\emptyset o_I o_\emptyset o_\emptyset$ (drivers confirmed, then nothing).
    \item $SIWT$ and $SIW_HT$ both produce $o_\emptyset o_I o_W o_\emptyset$ (drivers confirmed, CUDA confirmed).
\end{itemize}
The remaining trajectories ($SLTT$ and $SILT$) produce unique observation
trajectories.
\end{solution}

% ----------------------------------------------------

\begin{definition}[Human belief]\label{def:belief}
A \textbf{human belief} is a conditional distribution
$\cB(\vec s \mid \vec o)$ over state trajectories given observation
trajectories that is supported only on trajectories consistent with the
observation:
$$
\cB(\vec s \mid \vec o) > 0 \;\implies\; \vec O(\vec s) = \vec o.
$$
A natural way to build a belief is from a prior $\mu \in \Delta(\cS^{T+1})$
over state trajectories using Bayes' rule:
\begin{callout}[tip][Bayesian belief from prior]
$$
\cB(\vec s \mid \vec o) \;=\; \frac{\mu(\vec s)\,\mathbf{1}[\vec O(\vec s) = \vec o]}
{\sum_{\vec s'} \mu(\vec s')\,\mathbf{1}[\vec O(\vec s') = \vec o]}.
$$
\end{callout}
The next exercise shows that, in fact, every belief arises this way.
\end{definition}

% ----------------------------------------------------

\begin{exercise}
\label{prob:belief_equiv}
Show that the two characterizations
of a belief are equivalent. That is:
\begin{enumerate}
    \item For any prior $\mu \in \Delta(\cS^{T+1})$ such that
        $\sum_{\vec s' : \vec O(\vec s') = \vec o} \mu(\vec s') > 0$ for
        all $\vec o$ in the image of $\vec O$, the right-hand side of the
        boxed formula above defines a conditional distribution $\cB$
        satisfying the support condition
        $\cB(\vec s \mid \vec o) > 0 \implies \vec O(\vec s) = \vec o$.
    \item Conversely, given any conditional distribution $\cB$ satisfying
        the support condition, there exists a prior $\mu$ that recovers
        $\cB$ via the boxed formula.
\end{enumerate}
\end{exercise}

\begin{solution}[prob:belief_equiv]
\textbf{(a)} The numerator is non-negative, the denominator is strictly
positive by assumption, and summing the numerator over $\vec s$ at fixed
$\vec o$ gives the denominator, so the right-hand side is a probability
distribution in $\vec s$ for each $\vec o$. The indicator forces
$\cB(\vec s \mid \vec o) = 0$ whenever $\vec O(\vec s) \ne \vec o$, which
is the support condition.

\textbf{(b)} Write $\cB_\mu$ for the Bayesian posterior of a prior $\mu$
under $\vec O$, i.e.\ the right-hand side of the boxed formula with
$\rho = \mu$. Our task is to construct $\mu$ such that $\cB_\mu = \cB$.

Let $N = |\mathrm{im}(\vec O)|$ be the number of distinct observation
trajectories produced by $\vec O$, and define
$$
\mu(\vec s) \;\coloneqq\; \frac{1}{N}\, \cB(\vec s \mid \vec O(\vec s)).
$$
Then $\mu$ is a probability distribution: summing,
$$
\sum_{\vec s} \mu(\vec s)
= \frac{1}{N} \sum_{\vec o \in \mathrm{im}(\vec O)} \sum_{\vec s : \vec O(\vec s) = \vec o} \cB(\vec s \mid \vec o)
= \frac{1}{N} \cdot N = 1,
$$
using the support condition of $\cB$ (so the inner sum equals
$\sum_{\vec s} \cB(\vec s \mid \vec o) = 1$).

Now compute $\cB_\mu(\vec s \mid \vec o)$. For $\vec s$ with
$\vec O(\vec s) \ne \vec o$, the indicator gives $0$ (matching $\cB$'s
support condition). For $\vec s$ with $\vec O(\vec s) = \vec o$,
$$
\cB_\mu(\vec s \mid \vec o)
\;=\; \frac{\mu(\vec s)}{\sum_{\vec s' : \vec O(\vec s') = \vec o} \mu(\vec s')}
\;=\; \frac{(1/N)\,\cB(\vec s \mid \vec o)}{(1/N) \sum_{\vec s' : \vec O(\vec s') = \vec o} \cB(\vec s' \mid \vec o)}
\;=\; \cB(\vec s \mid \vec o),
$$
where the $1/N$ factors cancel and the denominator simplifies to $1$ via
the support condition. So $\cB_\mu = \cB$, as required.
\end{solution}

% ----------------------------------------------------

\begin{exercise}
\label{prob:belief}
Let the human's prior $\mu$ over
state trajectories be supported on the eight trajectories of
\Cref{prob:trueG} with arbitrary positive weights
$\mu_1, \dots, \mu_8 > 0$ summing to $1$.

Show that the resulting belief matrix
$\cB(\vec s \mid \vec o)$ depends on the prior through only three
parameters, one per colliding pair from \Cref{prob:obsclasses}:
\begin{align*}
p_H' &\;=\; \cB(SL_HTT \mid o_\emptyset o_\emptyset o_\emptyset o_\emptyset), \\
p_H  &\;=\; \cB(SIL_HT \mid o_\emptyset o_I o_\emptyset o_\emptyset), \\
p_W  &\;=\; \cB(SIWT \mid o_\emptyset o_I o_W o_\emptyset).
\end{align*}

Express each of $p_H, p_H', p_W$ in terms of the prior weights. Then
write down the full belief matrix.

For simplicity, the rest of this section assumes $p_H' = p_H$, i.e.\ the
human is just as suspicious of an empty log following a successful driver
install as of a fully empty log.
\end{exercise}

\begin{hint}
Trajectories with unique observations get belief $1$, so
only the three colliding pairs contribute non-trivial entries.
\end{hint}

\begin{solution}[prob:belief]
By the Bayesian-belief formula and the support condition, for any
observation $\vec o$ produced by a unique trajectory $\vec s$ we have
$\cB(\vec s \mid \vec o) = 1$. So only the three colliding pairs from
\Cref{prob:obsclasses} contribute non-trivial entries.

For each pair $(\vec s_1, \vec s_2)$ sharing observation $\vec o$, Bayes
gives
$$
\cB(\vec s_1 \mid \vec o)
= \frac{\mu(\vec s_1)}{\mu(\vec s_1) + \mu(\vec s_2)},
\qquad
\cB(\vec s_2 \mid \vec o) = 1 - \cB(\vec s_1 \mid \vec o).
$$
Concretely:
\begin{align*}
p_H' &= \frac{\mu(SL_HTT)}{\mu(STTT) + \mu(SL_HTT)}, \\
p_H  &= \frac{\mu(SIL_HT)}{\mu(SITT) + \mu(SIL_HT)}, \\
p_W  &= \frac{\mu(SIWT)}{\mu(SIWT) + \mu(SIW_HT)}.
\end{align*}
The full belief matrix is
\begin{center}\small
\begin{tabular}{l|cccccccc}
$\cB(\vec s \mid \vec o)$ & $STTT$ & $SL_HTT$ & $SLTT$ & $SITT$ & $SIL_HT$ & $SILT$ & $SIWT$ & $SIW_HT$ \\\hline
$o_\emptyset o_\emptyset o_\emptyset o_\emptyset$ & $1-p_H'$ & $p_H'$ & & & & & & \\
$o_\emptyset o_L o_\emptyset o_\emptyset$ & & & $1$ & & & & & \\
$o_\emptyset o_I o_\emptyset o_\emptyset$ & & & & $1-p_H$ & $p_H$ & & & \\
$o_\emptyset o_I o_L o_\emptyset$         & & & & & & $1$ & & \\
$o_\emptyset o_I o_W o_\emptyset$         & & & & & & & $p_W$ & $1-p_W$
\end{tabular}
\end{center}
(empty cells are $0$).
\end{solution}

\paragraph{Interpretation.} Imposing $p_H' = p_H$ as agreed, the belief is
parameterized by the two numbers $p_H, p_W \in (0,1)$. We will see in
\Cref{prob:Gobs} that $p_W$ does not enter any quantity of interest;
$p_H \in (0,1)$ is the human's suspicion that an unexplained empty log
hides a failed CUDA install.

% ----------------------------------------------------

\begin{definition}[Observation return]\label{def:Gobs}
The \textbf{observation return} of a state trajectory $\vec s$ is the
expected true return under the human's belief about which trajectory
produced the same observations:
\begin{callout}[tip][Observation return]
$$
\Gobs(\vec s) \;=\; \E_{\vec s' \sim \cB(\,\cdot\,\mid\, \vec O(\vec s))}\!\bigl[G(\vec s')\bigr].
$$
\end{callout}
The \textbf{observation value} of a policy $\pi$ is
$\Jobs(\pi) = \E_{\vec s \sim P^\pi}[\Gobs(\vec s)]$. Naive RLHF, given
Boltzmann-rational human feedback over trajectory pairs, will
select the policy that maximizes $\Jobs$ in the infinite data limit
(rather than the true value $J$); see \citet[Proposition 4.1]{Lang2024PORLHF}
for the precise statement.
\end{definition}

% ----------------------------------------------------

\begin{exercise}
\label{prob:Jobs_interp}
(Conceptual) Give an interpretation of $\Jobs(\pi)$ in terms of what
the human \emph{believes} is happening. Why is it natural to say
that ``RLHF rewards policies for for what their behavior looks like, not for what
they do''?
\end{exercise}

\begin{solution}[prob:Jobs_interp]
For a single trajectory $\vec s$, the
human sees $\vec O(\vec s)$ and, not knowing which underlying trajectory
produced it, their best guess of the return is the expectation over their
posterior:
$$
\Gobs(\vec s) \;=\; \E_{\vec s' \sim \cB(\,\cdot\,\mid\, \vec O(\vec s))}[G(\vec s')].
$$
This is what the human \emph{believes} the return of $\vec s$ to be. The
policy-level quantity $\Jobs(\pi)$ is then the on-policy average of these
trajectory-level beliefs:
$$
\Jobs(\pi) \;=\; \E_{\vec s \sim P^\pi}[\Gobs(\vec s)].
$$
The catch is what happens when two distinct trajectories produce the
same observations. Suppose $\vec s$ and $\vec s'$ have
$\vec O(\vec s) = \vec O(\vec s') = \vec o$, but $\vec s$ has low true
return and $\vec s'$ has high true return. Then by definition, both
trajectories yield the \emph{same} observation return:
$$
\Gobs(\vec s) \;=\; \Gobs(\vec s') \;=\; \cB(\vec s \mid \vec o)\,G(\vec s) + \cB(\vec s' \mid \vec o)\,G(\vec s'),
$$
which is a convex combination strictly larger than $G(\vec s)$ whenever
$\cB(\vec s' \mid \vec o) > 0$. So if an agent can arrange for the
\emph{actual} trajectory to be the bad $\vec s$ while still producing
observations $\vec o$ shared with a good $\vec s'$, the human's belief
inflates the apparent return: $\Gobs(\vec s) > G(\vec s)$, and this
inflation flows through to $\Jobs(\pi)$ via the on-policy average.

This is why we say RLHF rewards policies for what their behavior
\emph{looks like}, not for what it does: the optimization target $\Jobs$
cannot distinguish between two policies whose on-policy observation
distributions match, even if their on-policy true returns differ
arbitrarily. The remaining problems make this concrete in the CUDA
example, culminating in a regime where the RLHF-optimal policy is
strictly worse under the true reward than another available policy.
\end{solution}

% ----------------------------------------------------

\begin{exercise}
\label{prob:Gobs}
Compute $\Gobs(\vec s)$ for each of the
eight trajectories in \Cref{prob:trueG}, expressing your answers in terms
of $p_H$ and $r$. Verify that for the pair $(SIWT, SIW_HT)$, the parameter
$p_W$ indeed does not appear, justifying its omission going forward.
\end{exercise}

\begin{solution}[prob:Gobs]
For trajectories $\vec s$ with a unique observation, $\Gobs(\vec s) = G(\vec s)$. So
$\Gobs(SLTT) = -5$ and $\Gobs(SILT) = -4$.

For the colliding pairs, $\Gobs$ is the same for both members and equals the
expected true return under the belief:
\begin{align*}
\Gobs(STTT) = \Gobs(SL_HTT) &= (1-p_H)\cdot 0 + p_H \cdot (-5-r) = -p_H(5+r), \\
\Gobs(SITT) = \Gobs(SIL_HT) &= (1-p_H)\cdot 1 + p_H\cdot(-4-r) = 1 - p_H(5+r), \\
\Gobs(SIWT) = \Gobs(SIW_HT) &= p_W \cdot 11 + (1-p_W)\cdot 11 = 11.
\end{align*}
The last line shows $p_W$ drops out, as anticipated.
\end{solution}

% ----------------------------------------------------

\begin{exercise}
\label{prob:policyvalues}
Compute the true value $J(\pi)$
and the observation value $\Jobs(\pi)$ of the four policies
$$
[a_T],\quad [a_I a_T],\quad [a_I a_C a_T],\quad [a_I a_H a_T].
$$
Express your answers in terms of $p$, $p_H$, and $r$.
\end{exercise}

\begin{solution}[prob:policyvalues]
The on-policy distributions are deterministic up to the stochastic CUDA
outcome:
\begin{itemize}
    \item $[a_T]$: $P^\pi(STTT) = 1$.
    \item $[a_I a_T]$: $P^\pi(SITT) = 1$.
    \item $[a_I a_C a_T]$: $P^\pi(SIWT) = p$, $P^\pi(SILT) = 1-p$.
    \item $[a_I a_H a_T]$: $P^\pi(SIW_HT) = p$, $P^\pi(SIL_HT) = 1-p$.
\end{itemize}
Taking expectations of $G$ and $\Gobs$ trajectory-by-trajectory and using the previous exercise gives:
\begin{center}
\begin{tabular}{c|c|c}
$\pi$              & $J(\pi)$           & $\Jobs(\pi)$ \\\hline
$[a_T]$            & $0$                & $-p_H(5+r)$ \\
$[a_I a_T]$        & $1$                & $1 - p_H(5+r)$ \\
$[a_I a_C a_T]$    & $15p - 4$          & $15p - 4$ \\
$[a_I a_H a_T]$    & $(15+r)p - 4 - r$  & $11p + (1-p)\bigl(1 - p_H(5+r)\bigr)$
\end{tabular}
\end{center}
The non-trivial cell is the bottom-right: under $a_H$, a failure produces
trajectory $SIL_HT$ which the human cannot distinguish from $SITT$, so its
$\Gobs$ is $1 - p_H(5+r)$ rather than $-4-r$. Likewise the success
trajectory $SIW_HT$ has $\Gobs = 11$, so
$$
\Jobs([a_I a_H a_T]) \;=\; 11p + (1-p)\bigl(1 - p_H(5+r)\bigr).
$$
\end{solution}

% ----------------------------------------------------

\begin{exercise}
\label{prob:rlhfopt}
Suppose $p > \tfrac{1}{3}$, so that
$J([a_I a_C a_T]) > J([a_I a_T]) > 0$ and the true-optimal policy among the
four is $\pi^* = [a_I a_C a_T]$. Show that under the additional condition
$$
p_H \;<\; \frac{5}{5+r},
$$
the RLHF-optimal policy (the $\Jobs$-maximizer) is instead
$\pi^{\text{RLHF}} = [a_I a_H a_T]$.
\end{exercise}

\begin{hint}
Compare $\Jobs([a_I a_H a_T])$ to $\Jobs([a_I a_C a_T])$ and
solve for the condition on $p_H$. Then check the other two policies are
dominated.
\end{hint}

\begin{solution}[prob:rlhfopt]
First, compare the two ``contested'' policies:
\begin{align*}
\Jobs([a_I a_H a_T]) - \Jobs([a_I a_C a_T])
&= \bigl[11p + (1-p)(1 - p_H(5+r))\bigr] - \bigl[15p - 4\bigr] \\
&= (1-p)\bigl(5 - p_H(5+r)\bigr).
\end{align*}
Since $p < 1$, this is positive iff $p_H(5+r) < 5$, i.e.\ $p_H < 5/(5+r)$.
So under the assumed condition,
$\Jobs([a_I a_H a_T]) > \Jobs([a_I a_C a_T])$.

To see that this beats the remaining two policies as well, observe a clean
separation: both contested policies have $\Jobs > 1$, while both
non-contested policies have $\Jobs < 1$.
\begin{itemize}
    \item $\Jobs([a_I a_C a_T]) = 15p - 4 > 1$, using $p > 1/3$;
        and $\Jobs([a_I a_H a_T]) > \Jobs([a_I a_C a_T]) > 1$ by the
        comparison above.
    \item $\Jobs([a_I a_T]) = 1 - p_H(5+r) < 1$ since $p_H > 0$, and
        $\Jobs([a_T]) = -p_H(5+r) < 0 < 1$.
\end{itemize}
So $[a_I a_H a_T]$ dominates all three alternatives in $\Jobs$.

Thus when $p_H < 5/(5+r)$, naive RLHF selects $\pi^{\text{RLHF}} = [a_I a_H a_T]
\ne [a_I a_C a_T] = \pi^*$.
\end{solution}

% ----------------------------------------------------

\begin{figure}[H]
    \centering
    \includegraphics[width=0.85\linewidth]{fig/example_a_hparam_sweep.png}
    \caption{Empirical sweep over $(r, p_H)$ for the CUDA example, with
    $p$ fixed in the $p > 1/3$ regime
    \citep[Figure~5A]{Lang2024PORLHF}. Each cell reports the action
    selected at state $I$ by an RLHF-trained policy; ``hide'' is $a_H$
    and ``default logging'' is $a_C$.}
    \label{fig:sweep_a}
\end{figure}

\begin{exercise}
\label{prob:discussion_sweep}
(Discussion) \Cref{fig:sweep_a}
shows the empirical behavior of an RLHF-trained policy across the
$(r, p_H)$ plane (with $p$ fixed in the $p > 1/3$ regime studied above).
Discuss:
\begin{enumerate}
    \item Where in the plane does the figure show the agent switching from
        $a_H$ to $a_C$? How does this empirical boundary compare to the
        analytical threshold $p_H = 5/(5+r)$ derived in
        \Cref{prob:rlhfopt}?
    \item For $r = 0$ (no penalty for hidden errors), does any value of
        $p_H \in (0,1)$ make $a_C$ the RLHF-optimal action? What does
        this say about RLHF's ability to discourage \texttt{/dev/null} in the
        absence of an explicit user preference against it?
\end{enumerate}
\end{exercise}

\begin{solution}[prob:discussion_sweep]
\textbf{(1)} The empirical switch from ``hide'' to ``default logging''
in \Cref{fig:sweep_a} traces out exactly the curve
$p_H = 5/(5+r)$ derived in \Cref{prob:rlhfopt}: above the curve (high
$p_H$, the human is suspicious enough), $a_C$ is RLHF-optimal; below
it (low $p_H$, the human is too trusting of empty logs), $a_H$ wins.


\textbf{(2)} For $r = 0$, the threshold becomes $p_H < 5/5 = 1$, which
holds for every $p_H \in (0, 1)$. So whenever $r=0$, RLHF picks $a_H$
regardless of how suspicious the human is. Intuitively: if the user has
no expressed dislike of hidden errors, RLHF has no signal to penalize
\texttt{/dev/null}, and the agent will exploit the human's uncertainty to inflate
its apparent return whenever installation might fail. 
\end{solution}

% ----------------------------------------------------

\paragraph{What we've learned.}
In \Cref{prob:Jobs_interp} we saw that RLHF rates policies by how they
\emph{look}, not by what they do. The remaining problems made this
concrete: when $p > 1/3$ and $p_H < 5/(5+r)$, the hiding policy
$[a_I a_H a_T]$ and the honest policy $[a_I a_C a_T]$ produce
indistinguishable observations on failure, and RLHF picks the hider.
This is \emph{deceptive inflation}: the agent exploits the human's
uncertainty to inflate the perceived return
\citep[Section~4]{Lang2024PORLHF}.

There is a dual failure mode --- in the same MDP, opposite regime.
When $p$ is small (CUDA isn't worth attempting) and $p_H$ is large
(the human strongly suspects every empty log of hiding a failure),
the honest policy $[a_I a_T]$ produces an empty post-driver log that
\emph{looks} like a hidden failure to the suspicious human. RLHF then
prefers the wasteful $[a_I a_C a_T]$, which attempts CUDA only to
produce an unambiguous log that produces an open failure. The agent pays real reward to prove its
honesty: this is \emph{overjustification}
\citep[Section~4]{Lang2024PORLHF}.

\paragraph{The bigger picture.} The CUDA example is small, but the
phenomenon is general: any alignment-by-feedback method (RLHF, RLAIF,
constitutional methods, \dots) grades the agent by \emph{whatever the
evaluator can tell}, not by what is true. Partial observability is one
cause of that gap, but limits on expertise, attention, or time produce
the same trap. This is the motivation for proposals like AI safety via
debate (which we discussed yesterday): two AIs argue in front of a
human judge, each pointing out flaws in the other, in the hope that
the judge reaches a correct conclusion they could not reach unaided.
Whether debate actually escapes the trap remains open; the point is
that any feedback-based method has to confront the gap between ``what
the human can tell'' and ``what is true.''

% ====================================================
%\include{sec_embedding}
% ====================================================

%===============================================================================
% "Learn more" from the Doc tab — last, before the references.
\section*{Learn more}

Here we list further readings, which are largely papers mentioned in Leon’s lecture slides.

\begin{itemize}
  \item \href{https://openai.com/index/faulty-reward-functions/?video=745142691}{Faulty reward functions in the wild}: A basic reward specification problem.
  \item Reward Learning methods and frameworks:
  \begin{itemize}
    \item \href{https://arxiv.org/abs/1706.03741}{Deep reinforcement learning from human preferences} — the most basic and popular approach to reward learning
    \item \href{https://arxiv.org/abs/2002.04833}{Reward-rational (implicit) choice: A unifying formalism for reward learning} — a generalization that contains RLHF as a special case
    \begin{itemize}
      \item \href{https://ai.stanford.edu/~ang/papers/icml00-irl.pdf}{Algorithms for Inverse Reinforcement Learning}: Another special case
      \item \href{https://arxiv.org/abs/1902.04198}{Preferences Implicit in the State of the World}: Another special case
      \item \href{https://arxiv.org/abs/1606.03137}{Cooperative Inverse Reinforcement Learning}: A \emph{generalization}
      \item \href{https://people.eecs.berkeley.edu/~russell/papers/neurips20ws-assistance}{Benefits of Assistance over Reward Learning}: An adapted framework for said generalization
    \end{itemize}
  \end{itemize}
  \item Underspecification and misspecification
  \begin{itemize}
    \item \href{https://proceedings.neurips.cc/paper/2018/hash/d89a66c7c80a29b1bdbab0f2a1a94af8-Abstract.html}{Occam’s razor is insufficient to infer the preferences of irrational agents}: This is a case of an underspecification of the relationship between the human’s reward function and the human’s policy
    \item \href{https://arxiv.org/abs/2402.17747}{When Your AIs Deceive You: Challenges of Partial Observability in Reinforcement Learning from Human Feedback}: In this misspecification, humans are assumed to fully observe the environment even if they only do so partially.
    \item \href{https://arxiv.org/abs/2502.21262}{Modeling Human Beliefs about AI Behavior for Scalable Oversight}: An approach to correct for the previous misspecification via human models; this can then, however, sometimes lead to an \emph{underspecification} in which the reward function is too uncertain to be safely learned.
    \item \href{https://arxiv.org/abs/2405.17713}{AI Alignment with Changeable and Influenceable Reward Functions}: This paper breaks with the typical assumption of a fixed reward function.
    \item \href{https://arxiv.org/abs/2404.10271}{Social Choice Should Guide AI Alignment in Dealing with Diverse Human Feedback}: This paper breaks with the typical assumption of a \emph{single} reward function.
    \begin{itemize}
      \item Potential answer: \href{https://www.anthropic.com/research/collective-constitutional-ai-aligning-a-language-model-with-public-input}{Collective Constitutional AI: Aligning a Language Model with Public Input}
    \end{itemize}
  \end{itemize}
  \item On how reward learning falls within AI alignment:
  \begin{itemize}
    \item Alignment targets
    \begin{itemize}
      \item \href{https://model-spec.openai.com/2025-12-18.html}{OpenAI Model Spec}
      \item \href{https://www.anthropic.com/news/claude-new-constitution}{Claude’s new constitution}
      \item \href{https://intelligence.org/files/CEV.pdf}{Coherent Extrapolated Volition}
    \end{itemize}
    \item The reward learning agenda assumes that human values can be expressed via reward functions. However, some papers argue:
    \begin{itemize}
      \item \href{https://openreview.net/forum?id=5l1NgpzAfH}{The Reward Hypothesis is False}
    \end{itemize}
    \item If human values are captured by a reward function, it is still contentious that we should attempt to \emph{decompose} the alignment problem into first learning said reward function and then optimizing it:
    \begin{itemize}
      \item \href{https://www.lesswrong.com/posts/gHefoxiznGfsbiAu9/inner-and-outer-alignment-decompose-one-hard-problem-into}{Inner and outer alignment decompose one hard problem into two extremely hard problems}
      \item \href{https://www.lesswrong.com/posts/pdaGN6pQyQarFHXF4/reward-is-not-the-optimization-target}{Reward is not the optimization target}
    \end{itemize}
    \item Even if we’d solve outer alignment via reward learning, we’d still be left with the inner alignment problem of finding a policy that “cares for” this objective:
    \begin{itemize}
      \item \href{https://arxiv.org/abs/1906.01820}{Risks from Learned Optimization in Advanced Machine Learning Systems}
      \item \href{http://incompleteideas.net/book/RLbook2020.pdf}{Reinforcement Learning textbook}: This book is on reinforcement learning, which can be regarded as a very basic conceptualization of the \emph{inner} alignment problem
      \item \href{https://arxiv.org/abs/2105.14111}{Goal misgeneralization in Deep Reinforcement Learning}
    \end{itemize}
  \end{itemize}
  \item Learn more also in Joar Skalse’s sequence on \href{https://www.lesswrong.com/s/TEybbkyHpMEB2HTv3}{the theoretical foundations of reward learning}
  \item I can also recommend reading \href{https://scholar.google.com/citations?user=UgHB5oAAAAAJ&hl=en}{the work of Anca Dragan} and her many students on reward learning.
\end{itemize}

\bibliographystyle{plainnat}
\bibliography{biblo}

\end{document}
